% The world of map example ex-milner: the nine facts of ms_base in
% theories/examples/MilnerSched.v, without their shifts x+ < y (ms_L).
% Columns are field labels, the top row holds the points rho.k, the
% bottom row their duals. Arrow styles tell the agent that needs the fact.
\begin{figure}[tb]
  \centering
  \begin{tikzpicture}[x=1.95cm, y=1cm]
    \foreach \k/\pt/\tn/\bn in {%
        0/{\rho}/c_0/d_1, 1/{\fld{\rho}{1}}/a_1/b_1, 2/{\fld{\rho}{2}}/c_1/d_2,%
        3/{\fld{\rho}{3}}/a_2/b_2, 4/{\fld{\rho}{4}}/c_2/d_0, 5/{\fld{\rho}{5}}/a_0/b_0} {
      \node[ptnode] (t\k) at (\k, 1.3) {$\pt$};
      \node[ptnode] (b\k) at (\k, 0) {$\dual{\pt}$};
      \node[endname, above=0pt of t\k] {$\tn$};
      \node[endname, below=0pt of b\k] {$\bn$};
    }
    % A_1
    \draw[factA] (b0) -- (b1);
    \draw[factA] (b0) to[bend right=38] (b2);
    \draw[factA] (t1) -- (b2);
    % A_2
    \draw[factB] (b2) -- (b3);
    \draw[factB] (b2) to[bend right=38] (b4);
    \draw[factB] (t3) -- (b4);
    % A_0
    \draw[factC] (t0.north east) .. controls (1.0, 2.45) and (3.4, 2.45) .. (t4.north west);
    \draw[factC] (t0.north east) .. controls (0.8, 3.0) and (6.7, 3.0) .. (b5.east);
    \draw[factC] (b4) -- (b5);
    % legend
    \begin{scope}[shift={(0, -1.75)}]
      \draw[factA] (0, 0) -- (0.3, 0); \node[legend] at (0.35, 0) {needed by $A_1$};
      \draw[factB] (1.5, 0) -- (1.8, 0); \node[legend] at (1.85, 0) {needed by $A_2$};
      \draw[factC] (3, 0) -- (3.3, 0); \node[legend] at (3.35, 0) {needed by $A_0$};
    \end{scope}
  \end{tikzpicture}
  \Description{The twelve points of the scheduler in six columns, one per field label, with the points in the top row and their duals in the bottom row. Nine arrows for the facts all point from a lower field label to a higher one.}
  \caption{The world of the scheduler (\cref{ex:milner}). An arrow from $x$ to $y$ is a fact $x \prec y$ of the first creation, the grey name next to a point is the end at that point, and every fact also holds from $\pplus{x}$ (not drawn).}
  \label{fig:world-example}
\end{figure}
